\documentclass{article}
\usepackage[UTF8]{ctex}
\usepackage{amsmath,amssymb,amsfonts}
\usepackage{enumitem}
\usepackage{geometry}
\geometry{a4paper, margin=1in}
% === 错题本排版优化（不参与题目标识） ===
\setlength{\parindent}{0pt}
\setlength{\parskip}{0.35em plus 0.10em minus 0.08em}
\setlength{\emergencystretch}{3em}
\setlist{itemsep=0.25em, topsep=0.40em, parsep=0pt}
\allowdisplaybreaks[2]
% === 错题本排版优化结束 ===
\begin{document}
\section*{郑州大学2026年数学分析真题}
\section*{129．郑州大学2026年数学分析试卷}

\subsection{题目1}
1．（10 分）求极限 \(\lim _{x \rightarrow 0^{+}} \ln (2026 x) \ln \frac{\ln 2026 x}{\ln \frac{x}{2026}}\) ．

\subsection{题目2}
2．（10 分）设 \(f(x)=\int_{\sqrt{\frac{\pi}{3}}}^{\sqrt{x}} \frac{1}{1+\cot u^{2}} \mathrm{~d} u\) ，求积分 \(\int_{0}^{\frac{\pi}{3}} \frac{f(x)}{\sqrt{x}} \mathrm{~d} x\) ．

\subsection{题目3}
3．（15 分）求幂级数 \(\sum_{n=2}^{\infty} \frac{x^{n}}{n^{2}-1}\) 的收敛域及和函数，并求 \(\sum_{n=2}^{\infty} \frac{1}{\left(n^{2}-1\right) 2^{n}}\) 的和．

\subsection{题目4}
4．（10 分）已知曲线积分 \(\int_{L} \frac{1}{\varphi(x)+y^{2}}(x \mathrm{~d} y-y \mathrm{~d} x)=A\)（常数），其中 \(\varphi(x)\) 是可导函数且 \(\varphi(1)=1 . L\)是绕原点 \((0,0)\) 一周的任意正向闭曲线，试求 \(\varphi(x)\) 及 \(A\) 。

\subsection{题目5}
5．（10 分）设 \(z=z(x, y)\) 由方程 \(a x+b y+c z=F\left(x^{2}+y^{2}+z^{2}\right)\) 确定，其中 \(F(u)\) 为可导函数，\(a, b, c\)为常数．证明：函数 \(z=z(x, y)\) 满足方程

\[
(c y-b z) \frac{\partial z}{\partial x}+(a z-c x) \frac{\partial z}{\partial y}=b x-a y .
\]

\subsection{题目6}
6．（10 分）设函数 \(f(x)\) 在 \([0,1]\) 上具有二阶连续导数，证明：在 \((0,1)\) 内存在一点 \(\xi\) ，使得

\[
\left|f(0)+f(1)-2 f\left(\frac{1}{2}\right)\right| \leq \frac{1}{4}\left|f^{\prime \prime}(\xi)\right| .
\]

\subsection{题目7}
7．（15 分）设 \(x_{0}=\sqrt{6}, x_{n+1}=\sqrt{6+x_{n}}(n \geq 0)\) ．

\begin{itemize}
  \item[（1）] 证明：数列 \(\left\{x_{n}\right\}\) 收敛，并求 \(\lim _{n \rightarrow \infty} x_{n}\) ．
  \item[（2）] 判断级数 \(\sum_{n=1}^{\infty}(-1)^{n-1} \sqrt{3-x_{n}}\) 的敛散性，如果收敛，判断是绝对收敛还是条件收敛。
\end{itemize}

\subsection{题目8}
8．（15 分）设函数列 \(\left\{f_{n}(x)\right\}\) 在 \([a, b]\) 上一致收敛于 \(f(x)\) ．且每个函数 \(f_{n}(x)\) 在 \([a, b]\) 上连续．\(\left\{x_{n}\right\} \subset[a, b]\) 满足 \(\lim _{n \rightarrow \infty} x_{n}=x_{0}\) ，证明： \(\lim _{n \rightarrow \infty} f_{n}\left(x_{n}\right)=f\left(x_{0}\right)\) ．

\subsection{题目9}
9．（20 分）设函数 \(f(x)\) 在 \([a, b]\) 上可导，在 \((a, b)\) 内二阶可导．\(f(a)=f(b)=0, f_{+}^{\prime}(a)>0, f_{-}^{\prime}(b)>0\) ．

\begin{itemize}
  \item[（1）] 证明：存在 \(c \in(a, b)\) ，使得 \(f(c)=0\) ．
  \item[（2）] 证明：存在 \(\xi \in(a, b)\) ，使得 \(f^{\prime \prime}(\xi)=0\) ．
  \item[（3）] 证明：存在 \(\xi_{1}, \xi_{2} \in(a, b)\) ，使得 \(f^{\prime \prime}\left(\xi_{1}\right)<0, f^{\prime \prime}\left(\xi_{2}\right)>0\) ．
\end{itemize}

\subsection{题目10}
10．（15 分）证明级数 \(\sum_{n=1}^{\infty}(-1)^{n-1} \frac{x^{n}}{n\left(1+x^{n}\right)}\) 在开区间 \((0,1)\) 上一致收敛于 \(f(x)\) ，且 \(f(x)\) 具有连续导数．

\subsection{题目11}
11．（20 分）设 \(f(x)\) 在 \([0,+\infty)\) 上连续，常数 \(a, b>0\) ．

\begin{itemize}
  \item[（1）] 若 \(\lim _{x \rightarrow+\infty} f(x)=k\) ，证明： \(\int_{0}^{+\infty} \frac{f(a x)-f(b x)}{x} \mathrm{~d} x=[f(0)-k] \ln \frac{b}{a}\) ．
  \item[（2）] 计算 \(\int_{0}^{+\infty} \frac{\arctan \left(2 x^{2}\right)-\arctan \left(x^{2}\right)}{x} \mathrm{~d} x\) ．
\end{itemize}
\end{document}
